%%%%     Please only edit bits which say to edit them     %%%%
%%%% They usually start with %///% to show where they are %%%%



%%% START : DO NOT EDIT %%%%

% Document setup
\documentclass[14pt]{extarticle}

\usepackage{graphicx,geometry,hyperref,enumitem,nicefrac,extsizes,color,nameref,calc,textcomp,csquotes,setspace,titlesec}

\hypersetup{
    pdfborder={0 0 0}
}

% Important: This must come near the start of the setup file

\makeatletter
\newcommand*{\currentname}{\@currentlabelname}
\makeatother

\makeatletter
\def\normaljustify{%
  \let\\\@centercr\rightskip\z@skip \leftskip\z@skip%
  \parfillskip=0pt plus 1fil}
\makeatother

\setlength\parindent{0pt}
\setlength{\parskip}{\baselineskip}
\setcounter{secnumdepth}{2}
\setlength{\fboxsep}{0.4cm}
\setlength{\fboxrule}{0.04cm}

\definecolor{niceblue}{RGB}{166,218,255}
\definecolor{nicebluedark}{RGB}{0,104,179}
\definecolor{nicegreen}{RGB}{202,255,166}
\definecolor{nicegreendark}{RGB}{71,179,0}
\definecolor{nicered}{RGB}{255,166,166}
\definecolor{nicereddark}{RGB}{179,0,0}
\definecolor{niceyellow}{RGB}{255,254,166}
\definecolor{niceyellowdark}{RGB}{219,216,0}

\newcommand{\sectionbreak}{\clearpage}
\newcommand{\nosectionbreak}{
    \renewcommand{\sectionbreak}{%
    }
}

\newcommand{\actioncol}{niceblue}
\newcommand{\notecol}{nicered}
\newcommand{\questioncol}{niceyellow}
\newcommand{\goalcol}{nicegreen}

\newcommand{\group}[1]{
	\begin{center}
	\parbox[c]{\linewidth} {
		#1
	}
	\end{center}
}

\newcommand{\colourbox}[2]{
	\begin{center}
	\fcolorbox{#1dark}{#1}{
		\parbox[c]{.9\linewidth} {
			#2 
		}
	}
	\end{center}
}

\newcommand{\numberlist}[2][]{
	\begin{enumerate}[series=#1]
		#2
	\end{enumerate}
	\vspace{-\topsep}
}

\newcommand{\resumenumberlist}[2]{
	\begin{enumerate}[resume=#1]
		#2
	\end{enumerate}
	\vspace{-\topsep}
}

\newcommand{\itemlist}[1]{
	\begin{itemize}
		#1
	\end{itemize}
	\vspace{-\topsep}
}

\newcommand{\valuelist}[2]{
	\begin{description}[align=left,labelwidth=*,leftmargin=#1]
		#2
	\end{description}
}

\newcommand{\actions}[2][\currentname acs]{
	\colourbox{\actioncol}{ 
		{\small\color{\actioncol dark} Actions}
		\resumenumberlist{ #1 }{ #2 }
	}
}

\newcommand{\restartactions}[2][\currentname acs]{
	\colourbox{\actioncol}{ 
		{\small\color{\actioncol dark} Actions}
		\numberlist[#1]{ #2 } 
	}
}

\newcommand{\notes}[1]{
	\colourbox{\notecol}{ 
		{\small\color{\notecol dark} Notes}\par
		#1
	}
}

\newcommand{\questions}[2][questions]{
	\colourbox{\questioncol}{ 
		{\small\color{\questioncol dark} Questions}
		\resumenumberlist{ #1 }{ #2 } 
	}
}
\newcommand{\restartquestions}[2][questions]{
	\colourbox{\questioncol}{ 
		{\small\color{\questioncol dark} Questions}
		\numberlist[#1]{ #2 } 
	}
}
\newcommand{\goals}[1]{
	\colourbox{\goalcol}{ 
		{\small\color{\goalcol dark} Goals}\par
		#1 
	}
}

\newcommand{\boxesdescription}{
	\vspace{-\topsep}
	\begin{description}[align=left,labelwidth=*,leftmargin=3cm]
		\item [{\color{\actioncol dark}Actions}] Stuff for you to do. They are highlighted in blue.
		
		\item [{\color{\notecol dark}Notes}] Notes about important stuff you need to be aware of (and possibly remember!). They are highlighted in red.
		
		\item [{\color{\questioncol dark} Questions}] Questions you should try to answer. Sometimes you'll need to write things down; other times you'll need to build something in the game. They are highlighted in yellow.

		{\bfseries Ask a helper or the teacher to check your answers.}
		
		\item [{\color{\goalcol dark} Goals}] Stuff you should have completed at the end of each section. They are highlighted in green.
	\end{description}
}

\newcommand{\screenshot}[3][.9\linewidth-0.8cm]{
	\begin{center}
	\includegraphics[width=#1]{img/screenshots/#2}\par
	{\small #3}
	\end{center}
}

\newcommand{\image}[3][.9\linewidth-0.8cm]{
	\begin{center}
	\includegraphics[width=#1]{img/other/#2}\par
	{\small #3}
	\end{center}
}

\newcommand{\todo}[1]{
	\image[7cm]{todo}{#1}
}

\newcommand{\spacer}{
	\par\vspace{14pt}
}

\newcommand{\answer}[1]{

{\bfseries #1}
}

%%% END : DO NOT EDIT %%%%



%///% Workshop info - edit these values for your workshop
\def\workshoptitle{Minecraft Redstone}
\def\workshopsubtitle{Part 1 of 2: The Basics of Redstone}
\def\workshopauthor{Ed Nutting}



%%% START : DO NOT EDIT %%%%

\begin{document}

\sffamily

% Title page - do not edit
\begin{titlepage}

	\newgeometry{margin=2cm}
	
	\begin{figure}
	\centering
	\begin{minipage}{.5\textwidth}
		\centering
		\includegraphics[width=.6\linewidth]{img/Bristol-University-Logo}
	\end{minipage}%
	\begin{minipage}{.5\textwidth}
		\centering
		\includegraphics[width=.4\linewidth]{img/Digimakers}
	\end{minipage}
	
	\vspace{18pt}
	
	\end{figure}
	
	\centering
	\normalfont
	{\Large Merchant Venturers School of Engineering\par Outreach Programme\par}
	\vspace{2cm}
	\sffamily
	{\huge\bfseries\workshoptitle\par}
	\vspace{0.5cm}
	{\LARGE\bfseries\workshopsubtitle}
	
	\vfill
	
	{\normalsize Created by\par
	\large\slshape\workshopauthor}
	
	\vfill
	
	{\normalsize Organised by\par
	\large\href{mailto:Caroline.Higgins@bristol.ac.uk}{Caroline.Higgins@bristol.ac.uk}}

	\vfill

	{\large Published on \today\par}
\end{titlepage}

% Document setup - do not edit
\newgeometry{margin=1.5cm}

\renewcommand{\abstractname}{Notes to Teachers \& Helpers}
\begin{abstract}

\vspace{-\topsep}
    
\noindent
\vspace{-\topsep}
    
\setlist{leftmargin=0mm}
\itemlist {
    \setlength{\itemsep}{18pt}
    
%%% END : DO NOT EDIT %%%%
    
% Put notes to teachers and session helpers as a bullet point list in this section

%///% Begin editable content

    \item This workshop is intended to last 1 to 1\( \frac{1}{2} \) hours.
    
    \item This workshop is intended for ages 9\textsuperscript{+} (years 5\textsuperscript{+}).

    \item The content is intended to be learnt through self-directed individual or pair game play, using this worksheet as a guide.

    \item The learning platform is Minecraft, the popular block-based building game.

    \item There are a number of versions of the Minecraft game, not all of which are compatible with this workshop:
    \valuelist{7cm}{ 
    	\item [Minecraft for Windows or Mac] This version {\bfseries is compatible}.\par {\footnotesize This is the normal version downloadable from the Minecraft website.}
    	\item [Minecraft for RaspberryPi] This version {\bfseries is not compatible}.\par {\footnotesize This version does not include the required Redstone features.}
    	\item [Minecraft Education Edition] This version {\bfseries is compatible}.\par {\footnotesize You may wish to set this up with your class before we arrive to run the workshop.}
    }

    \item Students should already be comfortable playing Minecraft.\par This means they should be able to move easily, place and destroy blocks, use items, access the inventory (in Creative mode) and be familiar with the various block types available in the game.

    \item This workshop teaches the following skills:\par {\footnotesize Items marked with an asterisk are directly relatable to the National Curriculum.}
    \valuelist{1cm}{
    	\item [-] Placing, destroying and designing basic circuits using Redstone in Minecraft
    	\item [*] Basic logic equations
    	\item [*] Logic gates: NOT, OR, NOR, AND
    	\item [*] Principles of digital design: Combining logic gates
    }
    
% End editable content

%%% START : DO NOT EDIT %%%%

}

\end{abstract}

% Document setup - do not edit
\newgeometry{margin=1.5cm}


%%% END : DO NOT EDIT %%%%

% Sections from this point onwards are your choice, though Introduction, Conclusion (/Wrap-up/What we learnt) and Extra Resources sections are recommended.

%///% Begin editable content

\section{Introduction}

Hi! In this short workshop we're going to try to introduce some of the concepts that electronic engineers use every day to design everything from your mobile phone, to computers and laptops to the settings of a toaster.

Let's get started. Each section is made up of four parts:
\boxesdescription

We'll also write some information between parts and include plenty of screenshots to help you out.

\group { 
    \actions{
    	\item Open Minecraft
    	\item Log in
    	\item Go to Single Player
    }

    \screenshot{Empty-World-List}{The Minecraft Single Player World List}
    
    \notes{
    	Click \enquote{Create New World}
    }
}

\actions{
    \item Create a new Creative world with the following setup:
    \valuelist{6cm}{
    	\item [Game Mode] Creative
    	\item [World Type] Superflat
    	\item [Preset] Redstone Ready
    	\item [Generate Structures] ON
    	\item [Allow Cheats] ON
    }
}

\group {
    \screenshot{World-Creation-1}{Create New World (Stage 1)}

    \notes{
    	Click \enquote{More World Options...}
    }
}

\group {
    \screenshot{World-Creation-2}{Create New World (Stage 2)}

    \notes{
    	Set \enquote{World Type} to \enquote{Superflat} by repeat clicking it.\par
    	\spacer
    	Then click \enquote{Customize}
    }
}

\group {
    \screenshot{World-Creation-3}{Superflat Customisation}

    \notes{
    	Click \enquote{Presets}
    }
}

\group {
    \screenshot{World-Creation-4}{Select Redstone Ready preset (bottom of list)}

    \notes{
    	Scroll to the bottom and click \enquote{Redstone Ready}
    	\spacer
    	Finally, click \enquote{Use Preset} then \enquote{Done} then \enquote{Create New World}.
    	\spacer
    	Wait for the world to load.
    }
}

\goals{
That's it for the introduction - you should now have created your new world ready for Redstone building.\par
\spacer
You can build Redstone in any type of world, but Redstone Ready worlds make it much easier.
}

\questions{
    \item What kind of block is the Redstone ready world made from?
    \item How many blocks vertically downwards are there till you reach the bedrock?
}





\section{Placing and Powering Redstone}

Let's get started with Redstone. It looks like this:

\image[4cm]{Redstonedust}{Redstone Dust}

You can get it from the Redstone tab of the Inventory:

\screenshot{Inventory-Redstone-tab}{Redstone tab of the Inventory}

\actions{
    \item Open the inventory
    \item Take some Redstone dust
    \item Take a Redstone torch
    \item Take a Redstone lamp
    \item Take a Redstone repeater
}

\image[4cm]{Redstone-repeater}{Redstone Repeater}

\subsection{Powering a lamp}

We can place Redstone dust on the ground to form wires. Wires move Redstone power around. 

\screenshot{Four-placed-Redstone-dust}{Four bits of Redstone dust placed on the ground}

\actions {
    \item Place some Redstone dust on the ground
    \item Place more Redstone dust to form a line
    \item Place a Lamp at one end of the line (on the end, not next to it)
    \item Place a Redstone torch at the other end of the line
}

\notes {
    Redstone torches look similar to normal torches - don't use the wrong one!
}

\screenshot{Power-to-lamp}{Powered wire going into lamp}

\screenshot{Powered-and-unpowered-wire}{Unpowered (left) and powered (right) Redstone wires}

\goals {
    The torch should be supplying power to the wire. The power should be traveling down the wire into the lamp, so the lamp should light up.
}


\subsection{Boosting power}

Redstone power gets weaker the further it is from the source of power (i.e. the torch). A \enquote{Repeater} allows us to boost the power. 

\notes {
    Redstone power runs out after 15 blocks.
}

\screenshot{Power-Distance-in-Wires}{2 wires where the power runs out}

\notes {
    We can use a repeater to boost the power. A repeater is a power source, but it only transmits power when it is supplied with power!
}

\screenshot{Repeater}{A repeater (Right: input wire, Left: output wire)}

\notes {
    A repeater only accepts power in on one side and only outputs power on the other side.
}

\screenshot{Repeater-as-Diode}{A repeater acts like a diode}

\actions {
    \item Destroy the lamp you placed
    \item Extend the wire so it is more than 15 blocks long (but less than 30!)
    \item Place a lamp at the end of the new, longer wire
}

\goals {
    Notice how the Redstone power \enquote{runs out} after 15 blocks so the lamp doesn't switch on.
}

\screenshot{Redstone-power-runs-out}{Redstone power decreases with each block}

\actions {
    \item Find the block of Redstone dust where the power first runs out
    \item Destroy the dust at this point
    \item Place a repeater where the dust was
}

\notes {
    Make sure your repeater points in the right direction.
}

\screenshot{Redstone-power-boosted}{Redstone power is boosted by the repeater}

\goals {
    The repeater boosts the power. Your lamp should now be switched on as power now reaches the end of the wire.
}

\questions {
    \item Which side of a repeater accepts power in?
    \item Which side of a repeater outputs power?
    \item How far (counted in blocks) does power travel out of a repeater?
}


\goals {
    You should now know how to:
    \itemlist {
    	\item Get Redstone from the inventory
    	\item Place Redstone dust to form a wire
    	\item Power and unpower Redstone wires using Redstone torches
    	\item Boost Redstone power using a repeater
    }
}




\section{Redstone Torches}

Redstone torches can be placed on top of or on the side of blocks. 

\screenshot{Redstone-torches}{Some Redstone torches}


\subsection{Flow of power from a torch}

Redstone torches output power in all directions except diagonally and except to the block they are placed on.

\notes {
    Redstone torches {\bfseries don't} output power to {\bfseries either:}
    \itemlist {
    	\item the block they are hanging on the side of {\bfseries or,} 
    	\item the block they are placed on top of.
    }
}

\notes {
    Redstone torches do not output power diagonally.
}

\actions {
    \item Place a block (of sandstone or anything else) on the ground
    \item Place a lamp on top of the block
    \item Destroy the block underneath the lamp
    \item Using one torch at a time, place torches around, underneath, on top of and hanging on the lamp. Work out which do and which don't turn on the lamp.
    \item Place a lamp on the ground
    \item Using one torch at a time, place torches around the lamp (including diagonally). Work out which do and which don't turn on the lamp.
}

\screenshot{Lamp-powered-to-side}{Lamp powered by torch to the side of it}

\screenshot{Lamp-powered-underneath}{Lamp powered by torch underneath it}

\screenshot{Lamp-not-powered}{Lamp not powered by torches placed on it or placed diagonally from it}

\actions {
    \item Build a block on the ground
    \item Place a block on top of the block you just built
    \item Place a Redstone torch on the side of the upper block
    \item Place a Redstone wire on the ground directly underneath the torch
}

\notes {
    This last exercise can be confusing as Redstone dust looks like it's lying on the ground - it doesn't look like a full block, but it does occupy a full block's space in the game. Thus the torch is directly above the Redstone dust and so it powers the dust.
}

\goals {
    The wire should be powered as the torch outputs power downwards to the block underneath it, which is the Redstone dust block in the wire.
}

\goals {
    You should now know which directions power flows from a Redstone torch.
}


\subsection{Torch on a block}

A Redstone torch placed on a block has a special feature. If a block is powered (just like we powered the lamp earlier) then any Redstone torch placed on the block will switch off! 

\notes {
    Redstone torches on unpowered blocks switch on. Redstone torches on powered blocks switch off. 
}

\actions {
    \item Place a block on the ground
    \item Place a torch on the side or top of the block you just made
    \item Build a wire on the ground going into the block (just like the wire we built earlier going into the lamp)
    \item Build a torch at the end of the wire
    \item Destroy and remake the torch while looking at the torch on the block. Notice how the torch on the block switches off when power is put into the block.
}

\notes {
    Almost any block can be powered by passing power into it using a wire. Alternatively, you can power blocks by placing a torch underneath them or, for some blocks, next to them, as we saw with the lamp earlier.
}

\screenshot{NOT-gates}{Torches on unpowered and powered blocks}

\goals {
    You should have seen the torch on the block switch off and on as you powered and unpowered the block.
}

\questions {
    \item What would happen if you powered a block using a wire from a torch placed on the same block?
    \item When you power on the wire, is there a delay before the torch on the block switches off?
}

\goals {
    You should now be able to work out whether a torch on a block will be on by seeing if there is any power flowing into the block the torch is placed on.
}

\group {

    We can describe what is going here like this: 

    {    
    	\begin{spacing}{1.2}
    	\rmfamily\centering\LARGE\itshape
    	A torch on a powered block is not on. \\
    	A torch on an unpowered block is not off.
    	\end{spacing}
    }
    
    In other words, if we call the power (or lack of power) going into the block the \enquote{input}, and the power supplied by the torch the \enquote{output}, then \enquote{the input is not the output}.
    
}
    
\notes {
    A torch on a block is called a {\bfseries NOT gate}. 
    \itemlist {
    	\item Power going into block is called the \enquote{input}.
    	\item Power going out from the torch on the block is called the \enquote{output}.
    	\item If the torch is on, we say the output of the gate is on.
    	\item If the torch is off, we say the output of the gate is off.
    	\item For a NOT gate (which is a torch on a block), the output is the opposite of any inputs.
    	\item We call this a NOR gate when there is more than one input
    	\item We write the names of gates in CAPITAL LETTERS to distinguish them from ordinary English words.
    }
}

\subsection{NOT and NOR gates}

\actions {
    \item Build a NOT gate
    \item Build a second NOT gate, three blocks in front of the first one
    \item Connect the output of the first NOT gate to the input of the second using a wire
}

\screenshot{Chained-NOT-gates}{Two NOT gates chained together}

\questions {
    \item What happens when you power the input of the first NOT gate?
    \item Is two NOT gates in a row like this, the same as a repeater?
}

\actions {
    \item Build a NOT gate on the ground
    \item Build three wires - one going into each side of the block of the NOT gate
    
    This makes it a {\bfseries NOR gate}.
    \item Build a wire from the output of the NOR gate (the torch on the block) going into a lamp
    \item Power and unpower the three inputs, one at a time and then multiple at once. Observe what happens to the output (by looking at the lamp).
}

\todo{Screenshot of the NOR-gate case}

\questions {
    \item What combination of inputs makes the lamp switch on?
    \item What combinations of inputs makes the lamp switch off?
    \item If you wire all three inputs together, then put them into the block, is it the same thing? 
}

\notes {
    Any point where two wires join is called an {\bfseries OR gate} because the power coming out of the join is on if either of the inputs is on.
}

\notes {
    A NOR gate is an OR gate followed by a NOT gate, hence the name \enquote{NOR} meaning \enquote{NOT-OR}.
}

\goals {
    You should now be able to build NOT and NOR gates using torches on blocks and predict the output of a NOT or NOR gate given a set of inputs.
}

\questions {
    \item What is the output of a NOT gate if its input is on?
    \item What is the output of a NOR gate, if it has three inputs of which two are on and one is off?
}


\section{Using NOR gates}

NOR gates are the most basic building block for all electronic circuits - including Redstone circuits. Amazingly, we can combine NOR gates to form anything from a simple lock to an entire computer or mobile phone!

\notes {
    NOR gates for the building block of all circuits.
}

We can build other kinds of gate from a NOR gate. Let's suppose we wanted to build a locked door. Our locked door might require 4 levers to be put into the correct positions for the door to open. 

Let's call the four levers of our lock inputs A, B, C and D. We want our door to open if A {\itshape and} B {\itshape and} C {\itshape and not} D are switched on. It sounds like we're going to need an AND gate! 

\subsection{AND gates}

An AND gate is easy to make. We have two inputs. We NOT each of them separately, and then NOR the outputs of the NOT gates together. It looks like this:

\todo{Screenshot of AND gate}

\notes {
    AND gate: NOT the inputs, NOR the results together\par
    \spacer
    The output of an AND gate is on only if both of its inputs are on.
}

\actions {
    \item Build a NOT gate (placing the torch on top of the block not on the side)
    \item Build a NOT gate next to this (this time, place the torch on the side of the block not on the top)
    \item Build a third NOT gate, like the first - torch on top, next to the second gate
    \item Place a piece of Redstone dust on top of the NOT gate (i.e. the block) in the middle
    \item Build separate wires going into the back of the two NOT gates (i.e. the blocks) on either side. These are the input wires.
    \item Build a wire from the torch in the middle to a lamp. This is the output.
    \item Power both of the input wires - notice how the lamp switches on.
}

\notes {
    This can be tricky to build. Look at the screenshot above if you get stuck or it doesn't work!
}

\questions {
    \item Which inputs to an AND gate have to be on to make the output (i.e. the lamp) switch on?
    \item Which inputs to an AND gate can be off to make the output (i.e. the lamp) switch off?
}

\goals {
    You should now be able to build an AND gate and predict the output of an AND gate from a given set of inputs.
}


\subsection{Locked doors}

\todo{All of this section}


\section{Wrap-up}

We hope you enjoyed this workshop! This workshop also has a second part where we teach you how to build more complex circuits like a Minecart Wave Machine and a Locked Corridor. Ask your teacher about it!

\goals{
    Hmmm...
    \itemlist {
    	\item TODO
    }
}




\section{Extra Resources}

Here's a few extra resources to help you along with this worksheet and some stuff to try at home.

\itemlist {
    \item \href{http://www.minecraft.net}{Minecraft website : http://www.minecraft.net}
}

\end{document}